
%% Beginning of file == mru.tex ==
%% 
%% Created 05-10-2008
\documentclass[%
%%%%%%%%%%%%%%%%%%%%
%% Final o draft? %%
final,
%draft,
%%%%%%%%%%%%%%%%%%%%
%% Total o no toal?%
total,
%nototal,
%%%%%%%%%%%%%%%%%%%%
%% Color o BW?%%%%%%
%slideColor,
slideBW,
%%%%%%%%%%%%%%%%%%%%
%% colorBG o no? %%%
%nocolorBG,
colorBG,
%%%%%%%%%%%%%%%%%%%%
%% PS o PDF %%%%%%%%
%ps,
pdf,
%%%%%%%%%%%%%%%%%%%%
%% otra opcion.. %%%
accumulate,
%noaccumulate,
%%%%%%%%%%%%%%%%%%%%
%%%%%%%%%%%%%%%%%%%%
%% Background %%%%%%
%frames
%azure
%contemporain
%nuancegris
%troispoints,
%lignesbleues
%darkblue
%alienglow
%autumn
%%%%%%%%%%%%%%%%%%%%
%%%%%%%%%%%%%%%%%%%%
]{prosper}

%%%%%%%%%%%%%%%%%%%%%%%%%%%%%%%%%%%%%%%%%%%%%%%%%%%%%%%%%%%%
%% Paquetes.. %%%%%%%%%%%%%%%%%%%%%%%%%%%%%%%%%%%%%%%%%%%%%
\usepackage[latin1]{inputenc}
\usepackage{pstricks,pstricks-add,pst-node,pst-text,pst-3d}
\usepackage{amsmath}
\usepackage{cancel}
%% Mate
\usepackage{pst-func}
\usepackage{pst-math}

%%%%%%%%%%%%%%%%%%%%%%%%%%%%%%%%%%%%%
%% Seguridad.. %%%%%%%%%%%%%%%%%%%%%%
\hypersetup{
pdfpagemode={FullScreen},
pdftitle={MRU e1: ver1},
pdfauthor={JM~Ramirez},
pdfcreator={JM~Ramirez},
pdfproducer={JM~Ramirez},
pdfkeywords={Movimiento,Rectilineo,Uniforme}}
%%%%%%%%%%%%%%%%%%%%%%%%%%%%%%%%%%%%%
%%%%%%%%%%%%%%%%%%%%%%%%%%%%%%%%%%%%%

%% You can insert definicions here..

%% DEFINICIONES MIAS..
%\def\psPiFour{12.566371}
%\def\psPiTwo{6.283185}
\def\psPi{3.14159265}
%\def\psPiH{1.570796327}
%\newdimen\pstRadUnit
%\newdimen\pstRadUnitInv
%\pstRadUnit=1.047198cm % this is pi/3
%\pstRadUnitInv=0.95493cm % this is 3/pi


%% Libreria de formulas %%%%%%%%%%%%%%%%%%%%%%%%%%%%%%%%%%%%
%% Velocidad en el Mov. Rec. Uniforme %%%%%%%%%%%%%%
\def\vmru{
%% Eq. 1 %%%%%%%%%%%%%
\begin{equation}
v=\frac{x}{t}
\label{vmru}
\end{equation}
%%%%%%%%%%%%%%%%%%%%%%%%%%%%%%%%%%%%%%%%%%%%%%%%%%%%
}

%% Distancia en el Mov. Rec. Uniforme %%%%%%%%%%%%%%
\def\xmru{
%% Eq. 2 %%%%%%%%%%%%%
\begin{equation}
x=vt
\label{xmru}
\end{equation}
%%%%%%%%%%%%%%%%%%%%%%%%%%%%%%%%%%%%%%%%%%%%%%%%%%%%
}

%% Tiempo en el Mov. Rec. Uniforme %%%%%%%%%%%%%%%%%
\def\tmru{
%% Eq. 3 %%%%%%%%%%%%%
\begin{equation}
t=\frac{x}{v}
\label{tmru}
\end{equation}
%%%%%%%%%%%%%%%%%%%%%%%%%%%%%%%%%%%%%%%%%%%%%%%%%%%%
}

%% Libreria de Notas    %%%%%%%%%%%%%%%%%%%%%%%%%%%%%%%%%%%%
%%%%%%%%%%%%%%%%%%%%%%%%%%%%%%%%%%%%%%%%%%%%%%%%%%%%
%%%%%%%%%%% NU para la velocidad %%%%%%%%%%%%%%%%%%%
%%%%%%%%%%%%%%%%%%%%%%%%%%%%%%%%%%%%%%%%%%%%%%%%%%%%
\def\vNUa%
{
%% Dibuja una cajita!! %%%%%%%%%%%%%%%%%%%
%% el background de gris..
\psframebox[linewidth=2pt,framearc=.3,fillstyle=solid,
fillcolor=lightgray]{
\begin{tabular}{c}
%%%%%%%%%%%%%%%%%%%%%%%%%%%%%%%%%%%%%%%%%%
{\bf Notas de Unidades}\\
{\sc magnitud}: Velocidad [v].
La velocidad puede venir en: \\

$\frac{{\rm km}}{{\rm h}}$ ; 
$\frac{{\rm m}}{{\rm s}}$ ; 
$\frac{{\rm cm}}{{\rm s}}$. \\

$1~\frac{{\rm km}}{{\rm h}}=
\frac{1000}{3600}~\frac{{\rm m}}{{\rm s}}
=\cancelto{0.277}{\frac{1}{3.6}}\frac{{\rm m}}{{\rm s}}$ \\

$1~\frac{{\rm m}}{{\rm s}} = 100~\frac{{\rm cm}}{{\rm s}}$

%%%%%%%%%%%%%%%%%%%%%%%%%%%%%%%%%%%%%%%%%%
\end{tabular}
}
%%%%%%%%%%%%%%%%%%%%%%%%%%%%%%%%%%%%%%%%%%

%% fin definicion de vNUa
}
%%%%%%%%%%%%%%%%%%%%%%%%%%%%%%%%%%%%%%%%%%%%%%%%%%%%
%%%%%%%%%%%%%%%%%%%%%%%%%%%%%%%%%%%%%%%%%%%%%%%%%%%%
%%%%%%%%%%%%%%%%%%%%%%%%%%%%%%%%%%%%%%%%%%%%%%%%%%%%

%%%%%%%%%%%%%%%%%%%%%%%%%%%%%%%%%%%%%%%%%%%%%%%%%%%%%%%%%%%%

%% Titulo y autor .. obligatorio..
\title{MRU e1: ver1}
\subtitle{Practica... y las puertas se abren}
\author{ejjr}
\email{ejjr.clases@gmail.com}

%%%%%%%%%%%%%%%%%%%%%%%%%%%%%%%%%%%
%%% Comienzo del documento.. %%%%%%
%%%%%%%%%%%%%%%%%%%%%%%%%%%%%%%%%%%
\begin{document}
\maketitle
%%%%%%%%%%%%%%%%%%%%%%%%%%%%%%%%%%

%%%%%%% Para pegar el log.. %%%%%%%%%%%%%%%%%%%%%%%%%%%%%%%
\Logo(-1,.6){\includegraphics[width=12cm]{images/citlogogrey.eps}}
%%%%%%%%%%%%%%%%%%%%%%%%%%%%%%%%%%%%%%%%%%%%%%%%%%%%%%%%%%%

%%%%%%% Advertencia para no dejar que se copien %%%%%%%%%%%
\slideCaption{{\gray \large Esta nota tiene la intenci\'on de que \'este
material no sea entregado como tarea.}}
%%%%%%%%%%%%%%%%%%%%%%%%%%%%%%%%%%%%%%%%%%%%%%%%%%%%%%%%%%%

%________________________________________________________________(ult. col, 65)
% ========================== slide1 ============================= 
\begin{slide}{Movimiento Rectil\'ineo Uniforme}
% 
% -Paragraph1 
%----------------------------------------------------------------
%----------------------------------------------------------------
% -Idea1 
% ============================================================== 
 La f\'ormula principal:

                                                               
% ============================================================== 
% -Idea2 
% ============================================================== 
 
{\huge \vmru}
\vNUa
                                                               
% ============================================================== 
% -Idea3 
% ============================================================== 
                                                               
                                                               
                                                               
% ============================================================== 
%----------------------------------------------------------------
%----------------------------------------------------------------
\end{slide}
%________________________________________________________________(ult. col, 65)
% ========================== slide2 ============================= 
\begin{slide}{MRU}
% 
% -Paragraph1 
%----------------------------------------------------------------
%----------------------------------------------------------------
% -Idea1 
% ============================================================== 

Recu\'erdalo por sus unidades:

% ============================================================== 
% -Idea2 
% ============================================================== 

\begin{equation*}
10\frac{\text{km}}{\text{h}} =
\frac{\text{Se recorren 10 km ({\red distancia})}}
{\text{en 1 hora ({\red tiempo})}}
\end{equation*}
                                                               
% ============================================================== 
% -Idea3 
% ============================================================== 
 
Las otras dos f\'ormulas que se derivan de (\ref{vmru}) son:

{\large \xmru \\ \tmru}
                                                              
% ============================================================== 
%----------------------------------------------------------------
%----------------------------------------------------------------
\end{slide}
%________________________________________________________________(ult. col, 65)
% ========================== slide3 ============================= 
\begin{slide}{MRU}
% 
% -Paragraph1 
%----------------------------------------------------------------
%----------------------------------------------------------------
% -Idea1 
% ============================================================== 
                                                               
Es uniforme porque la velocidad no var\'ia en el tiempo; \\
(v es el mismo en cada unidad de tiempo!)
                                              
% ============================================================== 
% -Idea2 
% ============================================================== 
 
%%%%%%%%%%%%%%%%%%%%%%%%%%%%%
%% Para centrar el grafico %%
\begin{centering}
%%%%%%%%%%%%%%%%%%%%%%%%%%%%%
%%%%%%%%%%%%%%%%%%%%%%%%%%%%%

%% Propiedades del grafico.. %%%
%% el eje de X comienza aqui..
\def\iX{0.} %%
\def\iY{0.} %%
%% el eje de y comienza aqui..
\def\jX{8.5} %%
\def\jY{4.5} %%

%%%%%%%%%%%%%%%%%%%%%%%%%%%%%%%%%%%%%%%%%%%%%%%%%%%%%%%%%%%%%%
%% Este es el pspicture enviroment.. el mas sencillo %%%%%%%%%
\psset{unit=0.75}
\begin{pspicture}(-0.5,-3.5)(8.5,4.5)
%%%%%%%%%%%%%%%%(x1 , y1  )(x2 , y2)
\psaxes{->}(0,0)(\iX,\iY)(\jX,\jY)
\uput[-90](4.5,-1){t(s)}        %% Eje X
\uput[180](-1,3.5){v(cm/s)}     %% Eje Y

%% Un recta y=ctte .. %%%%%%%
\def\iVctte{2.5}
%% (x1 , y1  )(x2 , y2=y1)
\psline[linestyle=dashed](0,\iVctte)(8.5,\iVctte)
%%%%%%%%%%%%%%%%%%%%%%%%%%%%%%%%%%%%%%%%%%%%%%%%%%%%%
%%% El do de latex-pstricks!!! %%%%%%%%%%%%%%%%%%%%%%
%%% quiero poner unos punto rojos sobre la recta %%%%
\multido{\iA=1+1}{7}{%
\psdot[linecolor=red,dotscale=2.](\iA,\iVctte)}%
%\pscircle[fillstyle=solid](\iA,2.5){0.2}}%
%%%%%%%%%%%%%%%%%%%%%%%%%%%%%%%%%%%%%%%%%%%%%%%%%%%%%
%%%%%%%%%%%%%%%%%%%%%%%%%%%%%%%%%%%%%%%%%%%%%%%%%%%%%                                                              
 
%%%%%%%%%%%%%%%%%%%%%%%%%%%%%%%%%%%%%
\end{pspicture}
%%%%%%%%%%%%%%%%%%%%%%%%%%%%%%%%%%%%%%%%%%%%%%%%%%%%%%%%%%%%%%
%%%%%%%%%%%%%%%%%%%%%%%%%%%%%%%%%%%%%%%%%%%%%%%%%%%%%%%%%%%%%%

%%%%%%%%%%%%%%%%%%%%%%%%%%%%%%
%% fin centrado del grafico %%
\end{centering}
%%%%%%%%%%%%%%%%%%%%%%%%%%%%%%
%%%%%%%%%%%%%%%%%%%%%%%%%%%%%%
                                                              
                                                               
% ============================================================== 
% -Idea3 
% ============================================================== 
                                                               
                                                               
                                                               
% ============================================================== 
%----------------------------------------------------------------
%----------------------------------------------------------------
\end{slide}

%________________________________________________________________(ult. col, 65)
% ========================== slide1 ============================= 
\begin{slide}{MRU}
% 
% -Paragraph1 
%----------------------------------------------------------------
%----------------------------------------------------------------
% -Idea1 
% ============================================================== 
 
Un << auto >> { anda } a { 152 km/h }. Comienza su pase por un puente y { 205 segundos } despu\'es est\'a en el otro extremo, saliendo completamente (del puente) { 15 segundos } despu\'es. ?`Cu\'al es la { longitud } del puente y cu\'al es la longitud del << auto >>? %% Hay 0 elementos de distancia %% Hay 2 elementos de tiempo %% Hay 1 elementos de velocidad
                                                              
                                                              
% ============================================================== 
% -Idea2 
% ============================================================== 
 
\begin{itemize}                                                              

\item El 1er paso en la resoluci\'on de cualquier problema, es la
identificaci\'on de los elementos del problema.            
              
\end{itemize}                                                 
% ============================================================== 
% -Idea3 
% ============================================================== 
 
%% Dibuja una cajita!! %%%%%%%%%%%%%%%%%%%
%% el background de gris..
\psframebox[linewidth=2pt,framearc=.3,fillstyle=solid,
fillcolor=lightgray]{
\begin{tabular}{c}
%%%%%%%%%%%%%%%%%%%%%%%%%%%%%%%%%%%%%%%%%%

{\bf Elementos del Problema}\\

$x=?$	\\

$v=\cancelto{5}{\frac{10}{2}}
\frac{metros}{segundos}$ \\

$t=30~minutos$	\\

%%%%%%%%%%%%%%%%%%%%%%%%%%%%%%%%%%%%%%%%%%
\end{tabular}
}
%%%%%%%%%%%%%%%%%%%%%%%%%%%%%%%%%%%%%%%%%%
 
% Generated with LaTeXDraw 2.0.0
% Fri Aug 01 20:56:48 CEST 2008
% \usepackage[usenames,dvipsnames]{pstricks}
% \usepackage{epsfig}
% \usepackage{pst-grad} % For gradients
% \usepackage{pst-plot} % For axes
\scalebox{1} % Change this value to rescale the drawing.
{
\begin{pspicture}(0,-0.3528646)(3.3785417,0.3528646)
\usefont{T1}{ptm}{m}{n}
\rput(8.4,2){\psframebox[linewidth=0.04]{Qu\'e f\'ormula uso?}}
\end{pspicture}
}

% ============================================================== 
%----------------------------------------------------------------
%----------------------------------------------------------------
\end{slide}


%________________________________________________________________(ult. col, 65)
% ========================== slide2 ============================= 
\begin{slide}{MRU}
% 
% -Paragraph1 
%----------------------------------------------------------------
%----------------------------------------------------------------
% -Idea1 
% ============================================================== 
 
{\huge                
\xmru
}                                             
 
% ============================================================== 
% -Idea2 
% ============================================================== 
 
Pero.. siempre debemos trabajar en el mismo sistema de unidades.
Es decir, si hay un metros/segundos,
debemos transformar los 
minutos a segundos :)

%%%%%%%%%%%%%%%%%%%%%%%%%%%%%%%%%%%%%%%%%%%%%%%%%%%%%%%%%%%%%%%%
%%%%%%%%%%%%%%%%% Cajita de Texto sencilla %%%%%%%%%%%%%%%%%%%%%
\begin{center}
\psdblframebox[linewidth=1.5pt,linecolor=red]{%
\parbox[c]{6cm}{
1 minuto = 60 segundos \\
Entonces 30 minutos son \\
30 $\times$ 60 segundos \\
30 minutos =
1800 segundos.

}}
\end{center}
%%%%%%%%%%%%%%%%%%%%%%%%%%%%%%%%%%%%%%%%%%%%%%%%%%%%%%%%%%%%%%%%
%%%%%%%%%%%%%%%%%%%%%%%%%%%%%%%%%%%%%%%%%%%%%%%%%%%%%%%%%%%%%%%%
                           
% ============================================================== 
% -Idea3 
% ============================================================== 
 
Finalmente..                                                              
                                                               
                                                               
% ============================================================== 
%----------------------------------------------------------------
%----------------------------------------------------------------
\end{slide}


%________________________________________________________________(ult. col, 65)
% ========================== slide3 ============================= 
\begin{slide}{MRU}
% 
% -Paragraph1 
%----------------------------------------------------------------
%----------------------------------------------------------------
% -Idea1 
% ============================================================== 
{\large
\begin{equation}
x=
(5\frac{m}{\cancel{s}})
(1800~\cancel{s})=
9000~m
\end{equation}}

                               
% ============================================================== 
% -Idea2 
% ============================================================== 
 
La distancia que recorre el ferrocarril es
\fbox{9000 m}

                                                             
% ============================================================== 
% -Idea3 
% ============================================================== 
                                                               
                                                               
                                                               
% ============================================================== 
%----------------------------------------------------------------
%----------------------------------------------------------------
\end{slide}


%%%%%%%%%%%%%%%%%%%%%%%%%%%%%%%%%%%%%%%%%%%%%%%%%%%%%%
%%%%%%%%%%%%%%%  Tail  %%%%%%%%%%%%%%%%%%%%%%%%%%%%%%%
%%%%%%%%%%%%%%%%%%%%%%%%%%%%%%%%%%%%%%%%%%%%%%%%%%%%%%

\end{document}
%%
%% End of file == mru.tex ==
%%
%% Modified 05-10-2008 
